\providecommand{\toplevelprefix}{../..}  %
\documentclass[../../book-main_es.tex]{subfiles}

\begin{document}

\chapter{Problemas abiertos y líneas de investigación sobre inteligencia}
\label{ch:future}

\begin{quote}
    ``\textit{El estudio se llevará a cabo sobre la base de la conjetura de 
        que cada aspecto del aprendizaje o cualquier otra característica de la 
        inteligencia puede describirse, en principio, con tal precisión que sea 
        posible hacer que una máquina lo simule. Se intentará determinar cómo hacer 
        que las máquinas usen el lenguaje, formen abstracciones y conceptos, resuelvan 
        clases de problemas que hoy están reservados a los seres humanos 
    y mejoren a sí mismas.}''

    $~$\hfill -- Propuesta para el programa de IA de Dartmouth, 1956
\end{quote}
\vspace{5mm}

Este manuscrito ha establecido que el problema común y fundamental subyacente a la inteligencia consiste
en {\em aprender representaciones parsimoniosas y estructuradas de distribuciones de baja dimensión
en un espacio de alta dimensión}. Desde esta perspectiva, hemos introducido de manera sistemática los principios matemáticos y los mecanismos
computacionales que permiten aprender y desarrollar dichas representaciones a partir de los datos observados, entendidos como información y conocimiento. 
La capacidad de autónomamente buscar descubrir estructura y parsimonia en un mundo
aparentemente aleatorio es una característica fundamental de cualquier inteligencia, natural o artificial.
Creemos que los principios y mecanismos expuestos en este libro son unificadores y universales, 
aplicables tanto a los animales como a las máquinas.

Esperamos que este libro ayude a los lectores a aclarar por completo el misterio que rodea 
las prácticas modernas de las redes neuronales artificiales profundas, mediante el desarrollo 
de una comprensión rigurosa de sus funciones y roles desde la perspectiva del aprendizaje de representaciones estructuradas de 
distribuciones de baja dimensión a partir de datos de alta dimensión.
A partir de esta comprensión, se hacen evidentes las capacidades y limitaciones de los modelos 
y sistemas de IA existentes:
\begin{enumerate}
    \item Los modelos y sistemas existentes de redes neuronales profundas en bucle abierto no alcanzan la categoría de sistemas de memoria completos, 
      capaces de aprender y automejorarse en forma autónoma.
    \item Las implementaciones actuales del aprendizaje de representaciones siguen siendo primitivas 
      y basadas en fuerza bruta, lejos de ser óptimas en términos de arquitecturas de red 
      o mecanismos de aprendizaje.
    \item Los modelos de IA actuales únicamente aprenden distribuciones de datos y realizan 
      inferencia bayesiana inductiva, lo que se diferencia de la inteligencia humana de alto nivel.
\end{enumerate}

Un objetivo de este libro es ayudar a los lectores a establecer una comprensión objetiva y sistemática 
de las tecnologías actuales de la inteligencia artificial y a reconocer qué desafíos pendientes y 
por tanto oportunidades persisten para el avance futuro de la inteligencia artificial. En el último capítulo, 
presentamos algunos de nuestros puntos de vista y proyecciones para el futuro. 

\section{Hacia la inteligencia autónoma: ¿cerrar el círculo?}
De la práctica en el campo de la inteligencia artificial durante la última década ha quedado 
claro que, con suficientes datos y recursos computacionales, se puede construir un modelo 
lo suficientemente grande y preentrenarlo para aprender la distribución a priori de todos los 
datos, digamos, \(p(\vx)\). Teóricamente, un modelo de tal magnitud podría memorizar casi todo 
el conocimiento existente sobre el mundo que se codifica en datos, como imágenes bidimensionales, 
formas tridimensionales, movimientos dinámicos y lenguajes naturales. Tal como se discutió al 
inicio de este libro, dicho modelo cumple una función análoga a la del ADN, que la vida emplea 
para registrar y transmitir conocimientos sobre el mundo.

De este modo, el modelo y la distribución aprendidos pueden emplearse para generar nuevas 
muestras de datos extraídas de dicha distribución. Asimismo, se puede emplear el modelo para 
realizar inferencias (por ejemplo, completado, estimación y predicción) utilizando el 
conocimiento memorizado en diversas condiciones, digamos, mediante el muestreo de la 
distribución \textit{posterior} \(p(\vx\mid \vy)\) ante una nueva observación \(\vy\).
Estrictamente hablando, tal inferencia es de tipo inductivo o estadístico.

Ningún modelo preentrenado, sin importar su tamaño, garantiza que la distribución aprendida 
hasta el momento sea completamente exacta o exhaustiva. Si nuestras muestras \(\hat{\vx}_{t}\)
extraídas del \textit{prior} actual \(p_{t}(\vx)\) o nuestras estimaciones \(\hat{\vx}_{t}(\vy)\) 
basadas en el \textit{posterior} \(p_{t}(\vx\mid \vy)\) resultan inconsistentes con la verdad 
\(\vx\), desearíamos corregir las distribuciones aprendidas:
\begin{equation}
    p_{t}(\vx) \rightarrow p_{t+1}(\vx) \quad \text{or}\quad
    p_{t}(\vx\mid \vy) \rightarrow p_{t+1}(\vx\mid \vy),
\end{equation}
basado en el error \(\ve_{t} = \vx_{t} - \hat{\vx}_{t}\). 
Este proceso se denomina corrección del error basada en retroalimentación del error, un 
mecanismo ubicuo en la naturaleza para el aprendizaje continuo. Sin embargo, cualquier 
modelo de horizonte abierto carece, en si mismo, de un mecanismo para revisar o 
mejorar la distribución aprendida cuando esta resulta errónea o incompleta.
Mejorar los modelos actuales de inteligencia artificial depende aún en gran medida 
de la intervención humana: supervisión o refuerzo mediante experimentación, evaluación 
y selección. Podríamos llamar a este proceso «selección artificial» de modelos de gran escala, 
en contraste con la selección natural que impulsa la evolución de la vida.

Tal como se estudió anteriormente en este libro (especialmente en el 
\Cref{ch:consistent-representations}), los sistemas de bucle cerrado alinean una 
representación interna con las observaciones sensoriales del mundo exterior. Pueden 
mejorar continuamente la distribución aprendida internamente y su representación para 
lograr (progresivamente) la consistencia o la autocoherencia. El siguiente paso 
inmediato consiste en desarrollar sistemas de memoria que verdaderamente operen en bucle cerrado, 
ilustrados en la \Cref{fig:open-to-closed}, que aprendan de manera autónoma y continua 
distribuciones generales de datos y se mejoren basándose en la 
retroalimentación de errores. 

Por consiguiente, la transición desde los modelos de bucle abierto entrenados de extremo a extremo,
actualmente populares , hacia sistemas de bucle cerrado con aprendizaje continuo 
\begin{equation}
    \text{modelos \textbf{de final abierto}} \; \Longrightarrow \;
    \text{sistemas \textbf{de bucle cerrado}}
\end{equation}
es la clave para que las máquinas emulen realmente la manera en que el cerebro 
animal aprende y aplica conocimientos en un mundo abierto.
Creemos que 
\begin{quote}
    \begin{center}
        \textit{los modelos de final abierto son aptos para un mundo cerrado, por grande 
        que sea; \\ mientras que los sistemas de bucle cerrado son adecuados para un mundo abierto, por pequeño que sea.}
    \end{center}
\end{quote}
En consecuencia, la así llamada «Inteligencia Artificial General» (IAG o AGI) jamás podrá 
lograrse simplemente mediante un sistema diseñado para memorizar todo el conocimiento y datos 
existentes del mundo. Ningún conocimiento, sin importar su magnitud, resulta verdaderamente 
generalizable.\footnote{Según Sir Karl Popper, un influyente representante de la filosofía científica, ¡toda 
teoría y todo conocimiento científico es falsable!} No obstante, es precisamente la capacidad 
de mejorar la memoria o el conocimiento previo y adaptarse a cualquier nuevo entorno lo 
que resulta verdaderamente generalizable. Por lo tanto, una inteligencia verdaderamente 
auténtica resulta, en sí misma, generalizable siempre que se materialice 
correctamente.\footnote{En nuestra opinión, la IAG debería representar la «inteligencia artificial genuina».}
Esto constituye exactamente la esencia del programa de cibernética desarrollado por Norbert Wiener 
durante la década de 1940, que discutimos en las primeras páginas de este libro.

\begin{figure}[t]
    \centering
    \includegraphics[width=0.9\linewidth]{\toplevelprefix/chapters/chapter9/figs/open-to-closed.png}
    \caption{De una red profunda abierta a un sistema de bucle cerrado.}
    \label{fig:open-to-closed}
\end{figure}

\section{Hacia la inteligencia natural: ¿Más allá de la retropropagación?}
La práctica reciente en el campo de la inteligencia artificial ha llevado a muchos a pensar que 
se debe construir un solo modelo grande destinado a aprender la distribución de todos los datos y 
memorizar todo el conocimiento. Si bien esta aproximación es tecnológicamente posible, dicha solución 
probablemente esté lejos de ser necesaria ni de ser la más eficiente. Como sabemos del entrenamiento de redes 
neuronales profundas, el único método escalable conocido para entrenar tales redes a escala 
es la retropropagación (BP) \citep{Back-Prop}. Aunque la retropropagación (BP) proporciona un 
mecanismo para corregir errores mediante señales de gradiente propagadas hacia atrás a través del 
modelo entero, se trata, sin embargo, de un enfoque con bastante de fuerza bruta que difiere 
sustancialmente de cómo aprende la naturaleza: la BP a escala de todo el modelo es una opción que la naturaleza 
simplemente no puede permitirse debido a su alto costo en energía y que no puede implementar 
debido a limitaciones físicas adicionales tales como restricciones energéticas o la disponibilidad de datos.

De modo más general, no podemos comprender realmente la inteligencia sin entender también cómo se 
puede implementar en forma eficiente, en términos tanto de cómputo como de datos. Es decir, se requiere abordar la complejidad computacional 
y de datos inherentes a la realización de los mecanismos vinculados al logro de los objetivos de la inteligencia.
Históricamente, nuestra comprensión de la inteligencia (artificial) ha evolucionado a través de varias fases: 
desde la complejidad de Kolmogorov ---incomputable--- hacia la entropía de Shannon; desde la computabilidad 
de Turing hacia la posterior noción de tractabilidad computacional,\footnote{Decimos que un problema es tratable
    si permite un algoritmo cuya complejidad es polinómica en el tamaño
del problema.} y hasta el fuerte énfasis en la escalabilidad de los algoritmos
en la práctica moderna del análisis de datos de alta dimensión
\citep{Wright-Ma-2022} y más recientemente la inteligencia artificial basada en redes profundas. 
Este camino de iluminación puede resumirse de la siguiente manera:
\begin{equation}
    \text{\textbf{incomputable}} \;
    \Longrightarrow \; \text{\textbf{computable}} \;
    \Longrightarrow \; \text{\textbf{tratable}} \; \Longrightarrow \;
    \text{\textbf{escalable}}.
\end{equation}
En gran medida, el éxito y la popularidad del aprendizaje profundo moderno y
la retropropagación se deben precisamente a que han ofrecido una implementación
hasta cierto punto escalable con plataformas informáticas modernas (como las GPUs) para
procesar y compactar grandes volúmenes de datos de alta dimensión. Sin embargo, dicha
implementación sigue estando lejos del nivel de eficiencia computacional y de datos con el que la naturaleza
alcanza la inteligencia.

Aún queda un margen considerable para mejorar la eficiencia de la inteligencia artificial,
de modo que pueda emular la eficiencia de la inteligencia natural,
la cual debería ser varios órdenes de magnitud mejor que
las actuales implementaciones basadas en fuerza bruta. Para ello, necesitamos descubrir, a partir
de los principios que subyacen a la inteligencia, 
nuevas arquitecturas de optimización y mecanismos de aprendizaje que permitan
aprender distribuciones de datos bajo condiciones físicas naturales y
restricciones de recursos, similares a las que enfrentan los seres inteligentes en
la naturaleza; por ejemplo, con acceso a solo unas pocas muestras cada vez y sometidos a fuertes 
limitaciones en términos de energía lo que no les permite actualizar
todos los parámetros del modelo de manera simultánea (mediante BP).

El marco fundamentado y el enfoque que se exponen en este libro pueden servirnos de guía
para descubrir esas nuevas arquitecturas de optimización y mecanismos de aprendizaje. Estas nuevas
arquitecturas y mecanismos deberían permitir el aprendizaje continuo en línea
y deberían poder actualizarse mediante una optimización, altamente localizada y esparsa, 
hacia adelante o hacia atrás.\footnote{Para aprender una distribución, una
    instanciación sencilla de estas propiedades deseadas se da en el caso más simple del
    PCA, con el método de PCA en línea presentado en el
    \Cref{ch:consistent-representations}. Una vez aprendido el modelo lineal de la
    distribución, la tarea aún más sencilla de realizar predicciones en línea
    con observaciones ruidosas puede llevarse a cabo fácilmente mediante métodos del tipo de mínimos cuadrados,
    como el filtro de Wiener \cite{Wiener-1942,Wiener-1949}. En el caso de los sistemas dinámicos lineales, esto se convierte en el famoso
filtro de Kalman; véase \cite[Apéndice B]{MaY2003} o \citep{peng2025mathematics}.}
  
\begin{figure}[t]
    \centering
    \includegraphics[width=0.99\linewidth]{\toplevelprefix/chapters/chapter9/figs/loops.png}
    \caption{Una arquitectura hipotética de la corteza cerebral. La corteza
        es un sistema de autocodificación masivamente paralelo y distribuido
        que consiste en una jerarquía de autocodificadores de bucle cerrado. Estos
        extraen información de sus entradas y maximizan la
        ganancia de información de las representaciones que generan. Dichos procesos se llevan a cabo en múltiples
    niveles de la jerarquía, lo que da lugar a distintos grados de abstracción de la información.}
    \label{fig:loops}
\end{figure}
Tal como nos ha enseñado la neurociencia, la corteza de nuestro cerebro
está formada por decenas de miles de columnas corticales
\citep{Hawkins-2021}. Todas las columnas corticales tienen estructuras
físicas y ---por tanto--- probablemente funciones similares. Son altamente paralelas y están distribuidas,
aunque están interconectadas en forma esparsa. Por lo tanto, creemos que para desarrollar un
sistema más eficiente y escalable para lograr representaciones más estructuradas (una memoria) del mundo, debemos considerar
arquitecturas que emulen las de la corteza. La \Cref{fig:loops}
muestra una arquitectura hipotética de este tipo: un sistema masivamente distribuido y
jerárquico compuesto por numerosos módulos de autocodificación de bucle cerrado
que funcionan en gran medida en paralelo.\footnote{En la bibliografía existen elementos de este tipo de
    arquitecturas hipotéticas en diversas formas, como
    las redes de codificación predictiva \citep{millidge2022predictive}, que cuentan
con una estrategia de optimización que utiliza actualizaciones locales eficientes.}
Estos módulos aprenden a codificar diferentes modalidades sensoriales o diferentes
proyecciones de datos de cada modalidad sensorial. Nuestro análisis en la
\Cref{sec:measurement-self-consistency} del
\Cref{ch:conditional-inference} sugiere que dicha detección
y aprendizaje paralelos de una distribución de baja dimensión en un espacio de alta dimensión son teóricamente
posibles. Luego se pueden entrenar autocodificadores de nivel superior (con pérdida) basándose
en las salidas de los de nivel inferior para desarrollar «abstracciones» más esparsas ---incluso discretas--- y
de nivel más alto, tales como el lenguaje, que las representaciones aprendidas por los niveles inferiores.

La arquitectura de sistemas distribuida, jerárquica y de bucle cerrado ilustrada en 
la \Cref{fig:loops} comparte muchas de las características conocidas de la corteza cerebral. Dicha 
arquitectura podría abrir muchas más posibilidades que la actual, basada en una única arquitectura 
grande de bucle abierto (como el modelo Transformer). En particular, podría dar lugar a nuevos mecanismos de 
aprendizaje hacia adelante y hacia atrás, que resulten en una representación cada 
vez más estructurada de los datos sensoriales. Un sistema así 
emularía de verdad el funcionamiento de la memoria en nuestro cerebro. Además, nos permitiría 
transformar las actuales vías de escalado de la inteligencia artificial en un nuevo paradigma 
de aprendizaje continuo capaz de igualar o incluso superar la eficacia y la eficiencia de la 
inteligencia en la naturaleza:
\begin{equation}
    \text{\textbf{no computable}}
    \Longrightarrow \text{\textbf{computable}}
    \Longrightarrow \text{\textbf{tratable}} \Longrightarrow
    \text{\textbf{escalable}}
    \Longrightarrow
    \text{\textbf{natural}}.
\end{equation}

\section{Hacia la inteligencia científica: ¿más allá de la Prueba de Turing?}

\subsection{La evolución de la inteligencia en la naturaleza}
Como hemos comentado al inicio de este libro, en el \Cref{ch:intro},
la inteligencia en la naturaleza ha evolucionado a lo largo de múltiples fases y
se ha manifestado en al menos cuatro formas distintas:
\begin{equation}
    \text{\textbf{inteligencia filogenética}} \;
    \Longrightarrow \; \text{\textbf{ontogenética}} \;
    \Longrightarrow \; \text{\textbf{societaria}} \;
    \Longrightarrow \; \text{\textbf{científica}}.
\end{equation}
Todas las formas de inteligencia comparten el objetivo común (y la característica) de 
aprender información o conocimiento útil en forma de distribuciones de baja dimensión de 
datos de alta dimensión percibidos del mundo abierto. No obstante, difieren en los 
libros de códigos o esquemas de codificación específicos que adoptan para codificar 
la información o el conocimiento útil, en la información o el conocimiento específico 
aprendidos por el individuo o grupo inteligente, en los mecanismos computacionales de adquisición y mejora 
de la información o el conocimiento codificados y en el sustrato físico y la implementación 
de dichos mecanismos. De acuerdo con los conceptos y la terminología desarrollados en este 
libro, las cuatro etapas de la inteligencia se diferencian en los tres aspectos siguientes:
\begin{enumerate}
    \item El \textit{libro de códigos} y el esquema de codificación utilizado para codificar la
        información o el conocimiento deseados.
    \item La \textit{información} o los conocimientos adquiridos y codificados
        mediante el libro de códigos.
    \item Los \textit{mecanismos de optimización o aprendizaje} utilizados para crear y
        corregir la información o el conocimiento codificados.
\end{enumerate}
La \Cref{tab:four-stage-intelligence} resume sus principales características:
\begin{table}[h]
    \centering
    \resizebox{1\columnwidth}{!}
    {
        \begin{tabular}{| c | c | c | c | c |}
            \hline & \textbf{Filogenético} & \textbf{Ontogenético} &
            \textbf{Societal} & \textbf{Científico}\\
            \hline

            \textbf{Libro de código}  & ADNs & Neuronas/Cerebro & Lenguajes Naturales
            & Abstracciones matemáticas \\ [0.5ex]
            \hline

            \textbf{Información} & Genes & Memoria & Conocimiento
            empírico & Hechos científicos\\ [0.5ex]
            \hline
            \textbf{Optimización} & Refuerzo & Retroalimentación de error &
            Prueba y error & Teorizar y Falsar \\  [0.5ex]
            \hline
        \end{tabular}
    }
    \caption{Comparación de las cuatro etapas de la inteligencia en la naturaleza.}
    \label{tab:four-stage-intelligence}
\end{table}

Como ya sabemos, la humanidad ha logrado dos saltos cuánticos en materia de su inteligencia:
\begin{enumerate}
    \item El primero fue el desarrollo del lenguaje oral y escrito,
        lo que permitió a los seres humanos compartir y transmitir el conocimiento adquirido
        a lo largo de las generaciones, de manera muy similar a como lo hace el ADN en la naturaleza.
    \item El segundo fue el desarrollo de las matemáticas y la lógica formal
        que se convirtieron en el lenguaje de la ciencia moderna y que posibilitan la 
        metodología científica de la adquisición de conocimientos.
\end{enumerate}
El lenguaje de las matemáticas nos liberó de resumir el conocimiento a partir de
las observaciones únicamente en forma empírica y nos permitió formalizar
el conocimiento en forma de teorías verificables o falsificables mediante la deducción
matemática o la experimentación. A través de la formulación de hipótesis, la deducción
lógica y la verificación o falsificación experimental, ahora podemos
descubrir y desarrollar de manera proactiva nuevos conocimientos que van mucho más allá
de lo que se puede aprender a partir de las distribuciones de los datos
captados por sensores. Por ejemplo, las relaciones causales no pueden aprenderse
únicamente a partir de las distribuciones \citep{Pearl-2009}.

A pesar de las aparentes diferencias entre estas distintas etapas de
la inteligencia, la evolución de la inteligencia comparte una característica
común:\textit{todos ellos utilizan ciertas formas de
    retroalimentación\footnote{Pérdida ante el error y recompensa ante el éxito.} para crear
    y corregir la información aprendida, y la retroalimentación y la corrección se vuelven
    cada vez más frecuentes y eficientes a lo largo de la evolución: desde el refuerzo de bucle abierto hasta la autosupervisión de bucle cerrado, desde
el ensayo y error colectivo hasta la formulación proactiva de hipótesis y su refutación.}

\subsection{De la inteligencia inductiva a la deductiva}
Como se señaló en la introducción (\Cref{ch:intro}), el programa de
«inteligencia artificial» (IA) de 1956 tenía precisamente como objetivo estudiar cómo
dotar a las máquinas de inteligencia científica, es decir, de funciones de alto nivel
como la abstracción matemática, la inferencia causal, la deducción lógica y 
la resolución de problemas, que se cree que diferencian a los seres humanos de los animales:
\begin{equation}
    \text{inteligencia (animal) de \textbf{bajo nivel}} \; \Longrightarrow \;
    \text{inteligencia (humana) de \textbf{alto nivel}.}
\end{equation}
Como hemos aclarado en repetidas ocasiones en este libro, la mayoría de los avances tecnológicos
en inteligencia artificial durante aproximadamente la última década, aunque
se hayan llevado a cabo bajo el nombre de «AI», están en realidad más estrechamente relacionados
con el hecho de que las máquinas implementen mecanismos de aprendizaje de bajo nivel que comparten en gran medida
tanto los animales como los seres humanos, incluyendo la inteligencia filogenética,
ontogenética y societaria. En estos niveles, la inteligencia
crea información o conocimiento \textit{empírico} (a partir de distribuciones empíricas de datos) y lleva a cabo la inferencia (o predicciones)
que es, en su mayoría, de naturaleza \textit{inductiva} (casi) bayesiana. Los principios y métodos presentados en este libro tienen como objetivo revelar los objetivos matemáticos y
los mecanismos computacionales que subyacen a la inteligencia en estos niveles. Estos
proporcionan una sólida justificación teórica y evidencia computacional
para la siguiente afirmación:
\begin{quote}
    \centering
    \textit{Las máquinas pueden adquirir conocimientos empíricos y realizar
    inferencias inductivas.}
\end{quote}
Sin embargo, hasta ahora no hay pruebas rigurosas, ya sea teóricas o científicas, 
que indiquen que estos mecanismos por sí solos bastarían para alcanzar
el alto nivel de inteligencia humana que el programa original de «inteligencia artificial» realmente
pretendía comprender y quería que las máquinas imitaran o incluso superaran.

De hecho, al día de hoy, sabemos poco sobre cómo verificar
o certificar rigurosamente si un sistema es realmente capaz de alcanzar un nivel tan alto
de inteligencia, a pesar de que la Prueba de Turing, también conocida como el
Juego de la Imitación,\footnote{En la propuesta de Turing, el evaluador, o
    interrogador, es un ser humano. Sin embargo, la mayoría de los evaluadores humanos tienen una
    formación y unos conocimientos científicos limitados, y sus conclusiones pueden ser
subjetivas.} se propuso en 1950 \citep{Turing-1950}, como un 
intento inicial de abordar la pregunta fundamental:
\begin{quote}
    \centering
    \textit{¿Pueden pensar las máquinas?}
\end{quote}
Sin embargo, en la propuesta de Turing, el evaluador o interrogador es  
un ser humano. Pero la mayoría de los evaluadores humanos tienen un entrenamiento y unos conocimientos científicos limitados
y sus conclusiones pueden ser bastante subjetivas.

Durante mucho tiempo, no se consideró necesaria una definición precisa de dicha prueba, ya 
que las capacidades de las máquinas se hallaban muy por debajo de las de los seres humanos (o 
incluso de los animales). Sin embargo, a la luz de los recientes avances tecnológicos, muchos 
modelos y sistemas ahora pretenden haber alcanzado o incluso superado la inteligencia humana. Por lo 
tanto, ya va siendo hora de desarrollar una definición científica y ejecutable de la Prueba de Turing, 
es decir, un protocolo sistemático y riguroso para evaluar y certificar el nivel de inteligencia 
de un modelo o un sistema que pretenda ser «inteligente».

Como argumentó Turing en su artículo original, resulta bastante difícil
definir qué es «pensar» realmente. En su lugar, podríamos partir de algo que pueda considerarse más
tangible e intentar determinar primero si el sistema en cuestión es realmente capaz de
«comprender» ciertos conceptos o conocimiento abstracto. Por ejemplo, ¿cómo podemos
verificar si un modelo o sistema inteligente ha comprendido realmente
un concepto abstracto como: 
\begin{enumerate}
    \item la noción de igualdad entre dos cantidades,
    \item la noción de los números (naturales, racionales, reales o imaginarios),
    \item la noción de infinito y la inducción matemática,
    \item la noción de consistencia lógica\footnote{En matemáticas, si un conjunto de axiomas y sus resultados deducidos son consistentes y
        no se contradicen entre sí, entonces existen. Esto se conoce como el Principio de Hilbert. Esto difiere de la
        noción de «consistencia» del capítulo \ref{ch:consistent-representations}, que requiere que las representaciones aprendidas
        sean consistentes con las distribuciones de datos empíricos. Una noción de consistencia es de naturaleza inductiva
    y la otra es deductiva.} y la demostración por contradicción,
\end{enumerate}
por mencionar solo algunos? ¿O simplemente ha memorizado una enorme cantidad de
ejemplos de tales nociones? Cabe señalar que los modelos de lenguaje de última generación
siguen teniendo dificultades con preguntas matemáticas sencillas como: «¿Es
3,11 mayor o menor que 3,9?»\footnote{Algunos modelos han
    corregido sus respuestas a este tipo de preguntas mediante
    ingeniería dirigida, o han incorporado mecanismos de razonamiento adicionales
    que verifican las respuestas inmediatas y las corrigen durante la inferencia.
    Sin embargo, dejamos al lector la tarea de comprobar rigurosamente
    si algún modelo de lenguaje de vanguardia comprende realmente el
    concepto de los números (naturales, racionales, reales y complejos) y su
    aritmética asociada.}

A partir de este libro, sabemos que los sistemas modernos de IA generativa se basan principalmente 
en la inferencia bayesiana (y, por lo tanto, inductiva) para producir sus respuestas. ¿Pueden 
las máquinas también comprender verdaderamente la inducción matemática o la deducción lógica 
general y entender su diferencia con respecto a la inducción bayesiana? 
¿Cómo verificamos si un sistema comprende realmente las reglas de 
la deducción lógica y matemática y es capaz de aplicarlas con rigor? ¿O, 
una vez más, ha extraído y memorizado meramente los patrones estadísticos de tales deducciones 
a partir de una gran cantidad de instancias ---como los datos de cadena de pensamiento--- y 
simplemente ha explotado dichos patrones inductivamente durante la inferencia? Por lo tanto, la siguiente pregunta sigue abierta:
\begin{quote}
    \centering
    \textit{¿Hay alguna diferencia entre \textbf{memorizar} y
    \textbf{comprender}?}
\end{quote}
Solo aclarando esa diferencia podremos dar una
respuesta realmente sustantiva a la pregunta: \textit{«¿Pueden las máquinas comprender?»}

Probablemente sea más importante preguntarnos: ¿cómo verificamos si un sistema que se autodenomina
inteligente es capaz de \textit{crear conceptos abstractos o mejorar su propio conocimiento} de manera proactiva mediante métodos deductivos? Por ejemplo, ¿puede crear nuevos conceptos matemáticos
y descubrir nuevas leyes físicas o relaciones causales?
La conexión entre comprender y crear puede ser bastante profunda, como lo expresó de manera célebre el físico Richard Feynman:
\begin{quote}
    \centering
    \textit{«Lo que no puedo crear, no lo comprendo».}
\end{quote}
Esto de algún modo sugiere que la comprensión depende de si el conocimiento de uno se basa en haberlo creado a partir de ciertos principios fundamentales.

\subsection{Pruebas científicas de inteligencia}
Por lo tanto, ya es hora de que desarrollemos métodos de evaluación científicamente sólidos
para determinar a cuál de las siguientes categorías pertenece el nivel de
inteligencia de un sistema:
\begin{enumerate}[label={Nivel-\arabic*.}, leftmargin=*]
    \item el mero hecho de haber aprendido la distribución de baja dimensión de ciertos datos portadores de 
        información y haber sido capaz de regenerarlos;
    \item ser capaz de aprender en forma autónoma y continua
        nuevas informaciones o de corregir los conocimientos existentes a partir de nuevas
        observaciones y predicciones;
    \item ser verdaderamente capaz de comprender y crear 
        conceptos, reglas y leyes abstractas, así como saber cómo aplicarlos
        correctamente de manera deductiva;
    \item ser capaz de formular nuevas hipótesis científicas o conjeturas matemáticas y verificarlas o refutarlas mediante la experimentación o la deducción lógica.
\end{enumerate}

\begin{figure}[t]
    \centering
    \includegraphics[width=\linewidth]{\toplevelprefix/chapters/chapter9/figs/tests.png}
    \caption{Tres pruebas para diferentes niveles o tipos de
        inteligencia: la Prueba de Wiener para la inteligencia básica,
        la Prueba de Turing para la inteligencia a nivel humano y
    la Prueba de Popper para la inteligencia a nivel científico.}
    \label{fig:three-tests}
\end{figure}

Para evaluar y distinguir estos diferentes niveles de inteligencia,
sugerimos que probablemente deban existir al menos tres pruebas diferentes,
como se ilustra en la \Cref{fig:three-tests}:
\begin{enumerate}
    \item \textit{La Prueba de Norbert Wiener:} para determinar si un
        sistema es capaz de autocorregirse y generar nueva información o conocimientos 
        por sí mismo o si simplemente recibe y actualiza conocimiento
        de manera pasiva a través de refuerzos o supervisión externos;
    \item \textit{La Prueba de Alan Turing:} para determinar si un
        sistema es capaz de crear y comprender conceptos abstractos y
        conocimiento, o si simplemente aprende y memoriza estadísticas de sus
        símbolos y las utiliza para la inferencia bayesiana;
    \item \textit{La Prueba de Karl Popper:} para determinar si un
        sistema es capaz de explorar y de manera proactiva y de crear nuevo conocimiento abstracto
        mediante la formulación y verificación de hipótesis o teorías
        basadas en la consistencia experimental o la autoconsistencia lógica.\footnote{Con un espíritu similar, recientemente, Demis Hassabis, de Google 
        DeepMind ha sugerido una así llamada «Prueba de Einstein»: 
        Si un sistema inteligente, digamos, un modelo de lenguaje de gran tamaño, solo ha sido
        entrenado con hechos y conocimientos hasta el año 1905, ¿sería capaz
        de inventar o o descubrir la teoría de la relatividad especial?}
\end{enumerate}
Creemos que, para este tipo de pruebas, el evaluador o el interrogador
no debería ser una sola persona ni un grupo arbitrario de personas, sino más bien una
entidad con acreditación científica que siga un protocolo y un proceso
científicamente sólidos, de modo que las conclusiones sean rigurosas y fiables.

\subsection{De la compresión a la abstracción}
Como hemos visto a lo largo de este libro, la \textit{compresión} ha desempeñado
un papel fundamental en la ganancia de información o en la creación de conocimiento
empírico a partir de datos. Es el principio rector y un mecanismo universal para identificar una
distribución de datos de baja dimensión y organizar la información en ella codificada 
mediante una representación estructurada. En gran medida, explica la práctica actual de la
«inteligencia artificial» con redes profundas, que es en gran parte de naturaleza 
inductiva y bayesiana.\footnote{Los sistemas que poseen mecanismos deductivos integrados para 
llevar a cabo la verificación formal o la deducción lógica no entran en esta categoría.} Una cuestión pendiente para futuros
estudios es si \textit{la compresión por sí sola} es suficiente para alcanzar
toda la inteligencia deductiva de alto nivel mencionada anteriormente. 
En concreto, ¿son la abstracción matemática, la inferencia causal, la generación de hipótesis
y el razonamiento y la deducción lógicos algún tipo de formas ampliadas o
trascendidas de la compresión? ¿Existe alguna diferencia fundamental entre aprender distribuciones de datos 
a través de la compresión y formar conceptos y teorías de alto nivel mediante la abstracción (matemática)? De ahí que nos 
quede planteada otra cuestión abierta:
\begin{quote}
    \centering
    \textit{¿Hay alguna diferencia entre la \textbf{compresión} y
    la \textbf{abstracción}?}
\end{quote}
Solo conociendo la diferencia podremos intentar dar una respuesta significativa a la pregunta: 
«¿Pueden las máquinas crear conocimiento en forma proactiva mediante la generación de hipótesis y la deducción?»

En gran medida, la ciencia ---y su libro de código asociado, las matemáticas--- puede considerarse 
como la forma más avanzada y eficaz de inteligencia. Es en gran medida deductiva en su naturaleza, y exclusiva 
de seres humanos cultos e ilustrados. En cierta ocasión, el filósofo Sir Karl Popper sugirió:
\begin{quote}
    \centering
    \textit{«La ciencia puede describirse como el arte de la
    sistemática simplificación excesiva.»}
\end{quote}
¡Creemos que descubrir y comprender los principios matemáticos y los mecanismos computacionales subyacentes
a tal inteligencia de alto nivel, así como reproducirlos con éxito mediante máquinas,
será la última frontera para la ciencia, las matemáticas y la computación en su conjunto!

\end{document}
